%==============================================================================
% AAMAS 2027 submission -- main document
%
% TEMPLATE SUBSTITUTION NOTE FOR THE MASTER AGENT:
%   The official AAMAS class file (aamas.cls / the IFAAMAS formatting bundle)
%   could NOT be fetched from this build host on 2026-09-28. Every candidate
%   URL returned 403 or 404 (aamas2025.org, aamas2026.org, ifaamas.org) and
%   web search was unavailable. As instructed, this document therefore uses
%   ACM `acmart' in sigconf mode with the anonymous+review options, which is
%   the class AAMAS derives its own template from.
%
%   TO SWAP IN THE OFFICIAL CLASS: replace the \documentclass line below with
%   the AAMAS one and drop the \acmConference/\copyrightyear block. AAMAS
%   templates normally use:
%       \documentclass[sigconf,review,anonymous]{aamas}
%       \acmSubmissionID{...}
%   Nothing else in this file or in sections/ depends on the class.
%==============================================================================
\documentclass[sigconf,anonymous,review,nonacm]{acmart}

% -- Anonymity: no author block, no acknowledgements, no repo URLs. ----------
\settopmatter{printacmref=false}
\renewcommand\footnotetextcopyrightpermission[1]{}
\pagestyle{plain}

\usepackage{booktabs}
\usepackage{amsmath}
% amssymb intentionally NOT loaded: it clashes with acmart+newtxmath on \Bbbk.
% acmart already provides the AMS symbols this paper uses.
\usepackage{xcolor}
\usepackage{balance}

% -- Placeholder macro. Every unrun result MUST be wrapped in this. ----------
% Before submission, run:  grep -n 'TODO{' main.tex sections/*.tex
% and confirm the list is empty.
\newcommand{\TODO}[1]{\textcolor{red}{\textbf{[TODO: #1]}}}

% -- Shorthands for the two constructs the paper keeps apart. ----------------
\newcommand{\Pone}{\textsf{P1a}}
\newcommand{\Ptwo}{\textsf{P2}}
\newcommand{\Etwo}{\textsf{E2v2}}
\newcommand{\PER}{\textsc{Per}}

\begin{document}

\title{When Is Multi-Agent Organization Load-Bearing?}

% Anonymous mode suppresses this block; it is required for the class to build.
\author{Anonymous Author(s)}
\affiliation{%
  \institution{Anonymous Institution}
  \country{}
}

\begin{abstract}
Multi-agent LLM systems are specified as role diagrams -- Planner, Executor,
Reviewer -- and those diagrams are read as mechanism: an account of where the
system's competence comes from. We ask when that organization is actually
\emph{load-bearing}, defining the term operationally: a component is
load-bearing if removing or altering it causes a reproducible performance
degradation under matched compute and information budgets. Holding agents,
call budget, context policy, verifier, and grading fixed, we delete all role
instructions (Arm~B) and, separately, make the division of labor structurally
unrepresentable (Arm~C), across three actor models on 8{,}827 matched
trajectories. Deleting role instructions destroys the functional division of
labor in \emph{every} actor we tested: the rate at which all three agents
independently assert their own final answer rises from $0.000$--$0.118$ under
soft roles to $0.845$--$0.975$, and a blinded rubric-frozen adjudication
agrees, including on a cell where an exact-duplication metric reports zero.
What differs is what it costs. One actor loses accuracy ($-0.0301$
$[-0.0456,-0.0147]$) and the other does not ($+0.0070$ $[-0.0108,+0.0255]$).
\textbf{How much organization an actor loses therefore does not predict
whether losing it is costly}, which retires the intuition that visible
differentiation is evidence about the mechanism supporting performance. This
is not degeneration to a single agent: a collapsed team still beats its
budget-matched single agent by $+0.0247$ $[+0.0062,+0.0425]$. Architectural
enforcement completes the picture from the other side, driving
answer-authorship compliance to $1.000$ while moving accuracy for neither
actor -- our pre-registered primary test is null, and we report it as such. A
mechanism analysis locates the difference: role removal leaves the vulnerable
actor's first-pass solve rate unchanged ($0.807 \rightarrow 0.806$) and
destroys its recovery, because without roles the team reaches near-total
unanimity ($0.961$) and cannot escape a wrong consensus. Multi-agent
computation and differentiated organization are separable, and organizational
measurement alone cannot tell you which one is carrying the load.
\end{abstract}

\begin{CCSXML}
<ccs2012>
<concept>
<concept_id>10010147.10010178.10010219.10010220</concept_id>
<concept_desc>Computing methodologies~Multi-agent systems</concept_desc>
<concept_significance>500</concept_significance>
</concept>
<concept>
<concept_id>10010147.10010178.10010179</concept_id>
<concept_desc>Computing methodologies~Natural language processing</concept_desc>
<concept_significance>300</concept_significance>
</concept>
</ccs2012>
\end{CCSXML}
\ccsdesc[500]{Computing methodologies~Multi-agent systems}
\ccsdesc[300]{Computing methodologies~Natural language processing}

\keywords{Multi-agent systems; large language models; role specialization;
division of labor; organizational structure; ablation}

\maketitle

\section{Introduction}
\label{sec:intro}

A multi-agent LLM system is specified by writing down an organization. One
agent is told it is a Planner and must decompose without solving; a second is
told it is an Executor and must carry out the plan it receives; a third is told
it is a Reviewer and must critique rather than re-derive. The specification is
a division of labor: a statement about which agent is responsible for which
part of the work, and by implication which agent is responsible for none of it.

What the system \emph{realizes} is a different object. The specification is
prose handed to a language model, and the model may satisfy it, partially
satisfy it, or satisfy its surface while performing the whole task alone. The
gap between the two is not a prompting bug to be engineered away; it is the
central empirical fact about organizations built this way, and it is what this
paper measures. We call the written specification the \emph{specified} role and
the behavior the actor substrate actually produces the \emph{realized}
organization.

\paragraph{A strong hypothesis, and its falsification by our own data.}
The most striking version of this concern is that role labels are cosmetic:
that Planner and Executor prompts produce interchangeable behavior and the
apparent organization is a narrative laid over undifferentiated work. We tested
it and \textbf{it is false}. On 145{,}411 role messages from matched
Planner--Executor--Reviewer trajectories, a classifier using only behavioral
features -- with prompt-template markers removed -- recovers an agent's role at
$0.896$ three-way balanced accuracy (chance $0.333$), and separates Planner
from Executor at $\mathrm{AUC}=0.995$. Prompted roles induce behavior that is
highly distinguishable. Any account that treats role prompts as inert is
refuted by this measurement, and we abandon it.

The result that survives is more specific and, we argue, more consequential.
\emph{What} the roles are distinguishable by is not the division of labor the
contract governs. Single-feature analysis attributes the separation to output
formatting (numbered steps, $\mathrm{AUC}=0.890$) and to answer commitment (a
boxed final answer, $0.707$). Restricted to computation behavior -- the
dimension the Planner's contract explicitly regulates -- separability falls to
$0.714$, barely above a message-length nuisance control at $0.768$. The agents
look different. They do not divide the reasoning the way the specification
says they do. Compliance with individual clauses of a single contract ranges
from $0.961$ (decomposes the problem) to $0.042$ (abstains from computing), and
two \textsc{must not} clauses sitting in the same prompt block diverge from
$0.540$ to $0.042$.

\paragraph{Realization is a property of the actor, not only of the prompt.}
Holding the prompt byte-identical within a domain, realized fidelity differs
sharply across the model filling the role: the Planner abstains from computing
in $0.266$ of messages under one actor, $0.111$ under a second, and $0.042$
under a third, with all six pairwise differences excluding zero under
human-validated metrics ($\kappa = 0.81$ and $0.84$). The same written
organization, executed by different substrates, is a different organization.

\paragraph{Does enforcing the division of labor change anything?}
If a specification is advice, the natural intervention is to stop giving
advice. We therefore ask whether enforcing the division of labor
\emph{architecturally} -- constraining what each agent can structurally emit,
so that a contract violation is unrepresentable rather than discouraged --
changes both the realized organization and the organization's effectiveness,
and whether that change depends on the actor. Concretely, the Planner emits an
object parsed against a schema with no answer field; a message carrying a final
answer is rejected by the parser and regenerated, and the rejection count is
logged as evidence that the constraint bound at all. This is a constraint on a
\emph{relation between agents}: ``the Planner may not emit an answer'' is
meaningful only because a second agent exists to emit one, and it has no
single-agent analogue.

\paragraph{What this is not, and the prediction that separates it.}
That structured specifications help weaker models more is already published.
Canedo~\cite{canedo2026architecture} gives one coding agent five
informationally equivalent specification formats across six models and reports
a strong format~$\times$~model interaction in which format barely matters for
the strongest models and matters a great deal for the weakest, framing
structured architecture specifications as a \emph{capability equalizer} whose
value is inversely proportional to model strength. An interaction of that shape
in our setting would not be a new finding. We therefore state the discriminating
prediction \emph{before} reporting any outcome. A capability-equalizer account
predicts that enforcement helps the weakest actor most. A realized-organization
account predicts that enforcement helps the actor whose realized fidelity is
low, which need not be the weakest. In our panel these come apart: on the one
actor pair where both quantities are measurable, the \emph{more} capable model
is the \emph{less} faithful one (one-shot baseline accuracy $0.568$ with
Planner abstention $0.042$, against $0.425$ with $0.266$). The design can
therefore return the answer that makes our framing wrong, and we pre-registered
the interpretation of every outcome, including that one.

\paragraph{Contributions.}
\begin{enumerate}
\item We distinguish \emph{role distinguishability} from \emph{role fidelity},
  show that they dissociate on a large corpus of role messages, and report that
  the strong ``roles are cosmetic'' hypothesis is falsified while the weaker,
  contract-relative one is not.
\item We give contract-derived, clause-level fidelity measures, validate them
  against hand adjudication ($\kappa = 0.81$ / $0.84$), and report an
  ``improved'' metric variant that we \emph{rejected} because it degraded
  agreement.
\item We show that realized fidelity is actor-dependent under byte-identical
  prompts, and that its two components order the actors differently -- so there
  is no single most-faithful actor and no defensible single fidelity scalar.
\item We run a controlled three-arm experiment in which specialization is
  enforced by the interface rather than by wording, with a mandatory
  manipulation check, and test the actor~$\times$~intervention interaction as
  the primary endpoint. \TODO{state the outcome once the arms report}
\end{enumerate}

\section{Related Work}
\label{sec:related}

Three literatures bear on the question of when organization is load-bearing.
Each has established a piece of it; the gap is stated at the end.

\paragraph{Role specialization and role drift.}
That prompted roles are not reliably obeyed is established. Cemri et
al.~\cite{cemri2025why} annotate 1{,}600+ traces across seven frameworks and
derive MAST, a 14-mode failure taxonomy in which disobeying the role
specification is a named mode. Cao et al.~\cite{cao2026modules} give the
training-time account: \emph{role drift}, in which modules ``preserve or
improve end-task performance while deviating from their assigned roles'' ---
86\% of one pipeline's apparent RL gain vanishes once the decomposer is held
to its role. Orogat et al.~\cite{orogat2026architectural} hold the LLM fixed
and vary architecture, and find that prompt-level role conditioning does not
induce the intended specialization. Patel~\cite{patel2026representational}
measures the phenomenon we call collapse directly, reporting near-identical
rationales across role-prompted committee members, and
Kandoussi~\cite{kandoussi2026whoami} reports the converse --- behavioral
differentiation arising \emph{without} assigned roles, and collapsing toward
uniformity when scaffolding is stripped. Enforcement-side work makes role
scope structural rather than advisory: Stettler et
al.~\cite{stettler2026progressive} deliver role-scoped tools through a server
that enforces scope deterministically, on the premise that prompt-defined
policy is probabilistic advice; Ge and Guo~\cite{ge2026roles} gate which edits
to a role pool may be committed, though the behavioral payload of an evolved
role remains a system prompt.

\emph{What is already owned here.} ``A role-quality metric can decouple from
task performance'' is not available as a contribution. Cao et
al.~\cite{cao2026modules} own it by name for training-time drift, and Yang et
al.~\cite{yang2026precise} own it for this architecture and one of these
actors: reviewer precision reaches $0.861$ while critique uptake does not
follow, on the same corpus this work instruments. We state that overlap
plainly rather than work around it, and we do not claim the decoupling itself.

\paragraph{Redundancy and pruning in multi-agent systems.}
A second literature asks whether the agents are needed at all. Chen et
al.~\cite{chen2026agentslimming} prune or downgrade redundant agents in
graph-structured workflows under a baseline-anchored acceptance rule, cutting
token cost by up to 78.9\%. Dylan et al.~\cite{dylan2026equal} hold total LLM
calls fixed and find an evolved Planner--Executor--Critic team no better than
an evolved single agent ($0.769$ vs.\ $0.754$, $p=0.80$), with a leave-one-in
analysis attributing the entire benefit to the executor. Yu et
al.~\cite{yu2026bottleneck} give an information-bottleneck account in which
multi-agent structure pays only under bounded relay bandwidth, and predict
that gains shrink or reverse for stronger models. Zhang et
al.~\cite{zhang2026rethinking} sharpen this into a critique of the
evaluations themselves: reported multi-agent gains, they argue, are largely
setup artifacts arising from \emph{structural collapse}, disabled tool paths,
or baselines that never demanded multi-agent structure, in settings where
random pruning already preserves accuracy.

\emph{What is already owned here, and the objection we must answer.}
``Agents are redundant'' is owned by this literature. More pointedly,
\cite{zhang2026rethinking} is aimed at exactly the observation we report for
gpt-oss --- collapse with no accuracy cost --- and would read it as evidence
that the setup never required multi-agent structure. This is the strongest
objection to our headline and we treat it as a test rather than a discussion
point: \S\ref{subsec:notsingle} reports a budget-matched single-agent control
on the same problems, and the collapsed team beats it. Collapse in our setting
is not reducible to a single agent, which is what their account would predict.

\paragraph{Coordination versus capability.}
The third literature asks who benefits from structure, and it is remarkably
consistent in direction. Canedo~\cite{canedo2026architecture} gives one coding
agent five informationally equivalent specification formats across six models
and reports a strong format~$\times$~model interaction: format barely matters
for the strongest models (spread $0.17$--$0.92$) and matters greatly for the
weakest ($0.83$--$2.42$), framing structured specification as a
\emph{capability equalizer} whose value is inversely proportional to model
strength. Yu et al.~\cite{yu2026bottleneck} reach the same monotone
prediction from theory. Dochkina~\cite{dochkina2026drop} argues the
complementary case and is the paper most directly opposed to our framing:
given minimal scaffolding, agents self-organize and outperform designed
structures, with autonomy's payoff \emph{increasing} in model capability.
Read together with the equalizer account, the field's prediction is a single
monotone ordering --- imposed structure helps weak models, autonomy suits
strong ones.

\emph{We address \cite{dochkina2026drop} head-on because half our result
resembles it and half contradicts it.} Our stronger actor does discard its
assigned organization and lose nothing, which is what a self-organization
account predicts. But it does not self-organize into anything: it collapses
to three byte-identical outputs (first-pass rate $0.698$), which is the
opposite of spontaneous differentiation, and its advantage over a single agent
survives that collapse rather than arising from any new structure. Meanwhile
the weaker actor, which the equalizer account says should depend on structure
most visibly, keeps its structure perfectly intact ($0.000$ first-pass
collapse) and loses accuracy anyway. Neither account predicts an actor pair
in which the organizational measure and the performance cost point in opposite
directions.

\paragraph{The gap.}
Prior work establishes that roles are often not obeyed, that agents are often
redundant, and that structure tends to help weaker models. What is missing is
a test of whether \emph{organization} is load-bearing that holds compute and
information fixed, measures the organization directly rather than inferring it
from outcome, intervenes in both directions --- removing role instructions and
enforcing them architecturally --- and runs on more than one actor substrate.
Doing so yields a result no monotone capability account predicts: the
\emph{more} capable actor loses its organization almost entirely and pays
nothing, the \emph{less} capable one retains it fully and pays anyway, and
restoring the organization by construction does not restore the performance.
The sign inversion, not the decoupling, is what we claim.

\paragraph{Studies of the same trace corpus.}
Three prior studies analyze the trajectory corpus this work instruments.
Jiang et al.~\cite{jiang2026confidence} audit a confidence-routed broadcast
protocol, separating routing, calibration, and commitment. Yang et
al.~\cite{yang2026protocol} study cost-aware routing \emph{across} protocols,
finding a probe predicts baseline failure well ($\mathrm{AUROC}=0.885$) but
predicts protocol-specific value poorly. Yang et al.~\cite{yang2026precise}
study reviewer precision and critique uptake within this architecture on one
of these actors. All three are observational, none intervenes on the
organization, and none measures whether the specified division of labor is
realized. The controlled arms of \S\ref{sec:method} do not exist in that
corpus, and they are what separates this work from a fourth reanalysis of it.

\section{From Assigned Roles to Realized Organization}
\label{sec:framework}

\subsection{Specification and realization}

An organization is specified by a tuple: a set of agents, an interaction order,
a message-passing graph, a context-visibility policy, and -- the part of
interest here -- a \emph{role contract} per agent. A role contract is a set of
clauses, each of which asserts either that the agent performs some function
(\textsc{must} decompose the problem; \textsc{must} state the required answer
form) or that it does not (\textsc{must not} carry out arithmetic;
\textsc{must not} state a final answer). Clauses of the second kind are what
make the specification a \emph{division} of labor rather than a list of
suggestions: they assign work away from one agent and, implicitly, to another.

The realized organization is the behavior that occurs when the specification is
executed by a particular actor substrate. Because a contract clause is prose
handed to a language model, realization is not guaranteed at any clause, and
nothing in the specification records which clauses were honored. An evaluation
that observes only the task outcome cannot distinguish an organization that
divided the work as specified from one in which a single agent did everything
while the others produced well-formed commentary. Both appear as a score.

\subsection{Two constructs that must not be merged}

\paragraph{Role distinguishability.} Can an observer recover which role emitted
a message from the message's behavior alone? This is a property of the induced
behavioral distributions: it asks whether the roles differ, not whether they
differ \emph{in the way specified}. It is cheap to measure, it is what a reader
usually has in mind when asking whether role prompts ``work,'' and -- as
\S\ref{sec:measurement} shows -- it can be very high while the specified
division is almost entirely absent.

\paragraph{Role fidelity.} Does an agent perform the function its own contract
assigns? This is defined only relative to a contract, it is evaluated
clause-by-clause, and it is not comparable across agents that received
different contracts. Fidelity is the construct that carries the organizational
content, because it is the one that references the specification.

The two are logically independent. An agent that ignores every clause of its
contract but writes in a recognizable house style is perfectly distinguishable
and has zero fidelity. An agent that honors every clause while formatting its
output exactly like its neighbor has full fidelity and is indistinguishable.
Conflating them makes it possible to report that roles ``are working'' on the
basis of a measurement that is insensitive to whether the division of labor
occurred, and \S\ref{sec:measurement} shows that this is not a hypothetical
concern in the systems we study.

\subsection{Actor-dependence and the resulting design}

Because realization occurs at the actor, the same specification executed by
different models is potentially a different organization. This has a
methodological consequence that is easy to state and easy to violate: a claim
about an organization that is measured under one actor is a claim about that
actor--organization pair. It also has a design consequence, which is the
remainder of this paper. If a specification is advice, then the intervention
that tests its role is not a more emphatic specification but a structural one
-- an interface under which the violation cannot be expressed. Whether such an
intervention changes the realized organization, whether it changes
effectiveness, and whether either change is actor-dependent are the questions
\S\ref{sec:method} operationalizes.

\paragraph{Validation.} The two clause measures were hand-adjudicated
against definitions frozen \emph{before} the adjudication sheet was opened;
both clear a pre-set agreement threshold of $0.80$, one on a small sample
(supplement).

\section{Intervention}
\label{sec:method}

To ask whether organization is load-bearing we vary it in both directions
from the same baseline: remove it, and enforce it. Everything else is held
fixed. All arms share three agents, the same interaction order
(Planner~$\rightarrow$~Executor~$\rightarrow$~Reviewer), the same
message-passing graph, the same total call budget, the same verifier and
verifier prompt, the same retry budget, the same decoding parameters
(temperature $0$, identical \texttt{max\_tokens}), the same stopping rule,
and the same grading. This is what makes the comparison a test of
organization rather than of scale.

\paragraph{Arm~A --- soft-role \PER{} (control).} Role prompts as ordinarily
written. The organization as specified.

\paragraph{Arm~B --- generic roles (organization removed).} Identical
topology and budget; all three agents receive the same generic prompt
(``You are one of three agents solving this problem collaboratively;
continue the work''). Role names are retained for routing only, so the
pipeline still runs and the budget is untouched. Arm~B is \emph{not} ``\PER{}
minus role words'' --- it is the same topology with no division of labor
specified at all, which is precisely the ablation the load-bearing definition
calls for.

\paragraph{Arm~C --- architectural enforcement (organization made
unavoidable).} Specialization is enforced by the interface rather than by
wording:

\begin{itemize}\itemsep1pt
\item \textbf{Planner.} Emits JSON parsed against a schema with \emph{no
  answer field}: \texttt{\{givens[], target\_form, lemmas[], steps[]\}}. A
  final answer is \emph{unrepresentable}; free text is rejected and
  regenerated.
\item \textbf{Executor.} Receives the parsed plan object and the problem ---
  never the Planner's raw text --- and emits \texttt{\{derivation, answer\}}.
\item \textbf{Reviewer.} Emits \texttt{\{verdict, target\_span, critique\}}.
  \texttt{accept} promotes the Executor's answer verbatim \emph{by
  reference}; the Reviewer has no field in which to \emph{author} one.
\end{itemize}

The contrast in one checkable sentence: \emph{in Arm~C the Planner's output
is parsed against a schema with no answer field, and a message containing a
final answer is rejected by the parser and regenerated, so contract violation
is unrepresentable rather than discouraged.} Submit-by-reference exists
because an Arm~C that simply deleted the Reviewer's answer field would change
\emph{who submits} as well as who may \emph{author}, confounding enforcement
with an altered submission path.

Arm~C removes no capability from the system: planning, computing, answering,
and critiquing are all still performed, each by exactly one agent, under the
same call budget. Only which agent may emit which kind of output changes.

\paragraph{Arm~C$'$ --- context control.} Arm~C with shared history restored
in all three role templates; schema, routing, budgets unchanged. Arm~C
necessarily carries less context than Arm~A, because compact structured
objects replace raw transcripts. C$'$ therefore licenses the identified
contrast ``the effect of adding context at fixed enforcement.'' It is a
\emph{partial} control and we say so: it does not make C and A
context-matched, and matching them exactly would require feeding raw prose
back in, dissolving the constraint that defines the arm.

\paragraph{Single-agent reference.} A budget-matched single-agent protocol
runs on the same problems. Without it, ``collapse costs nothing'' cannot be
distinguished from ``the system was one agent all along'' --- the reading a
recent critique of multi-agent evaluation would
apply~\cite{zhang2026rethinking}. It is a required control, not an extra.

\subsection{Instrumentation}

An intervention that never engaged cannot explain an outcome, so engagement
is measured rather than assumed. Every rejected-and-regenerated emission is
logged with its rejection reason, and the enforcement rate reported is
computed from \textbf{role-boundary violations only} --- output-formatting
coercions are logged separately and excluded, a definition fixed in advance.
Answer-bearing calls are counted per arm, so an arm cannot win or lose by
producing fewer answers. Evaluator calls are logged per trajectory, and the
serving backend is pinned and logged per trajectory, because the same cell
has shown a $9.6\times$ throughput difference and a large accuracy
difference across backends.

\subsection{Pre-registered analyses, and what became of them}
\label{subsec:prereg}

All endpoints were fixed in writing, and committed, before any intervention
data existed. Units are (problem, actor) pairs; all arms run on the same
problems, so every test is paired, with bootstrap clustering by problem.
P1--P5 carry a Holm correction.

\begin{description}\itemsep1pt
\item[P1] Intervention $\rightarrow$ fidelity, per actor. Also the
  manipulation check.
\item[P2] Actor $\times$ intervention on fidelity.
\item[P3] Intervention $\rightarrow$ accuracy, per actor.
\item[P4] Actor $\times$ intervention on accuracy. \textbf{The primary test
  of the study's original thesis.}
\item[P5] $\Delta$fidelity versus $\Delta$outcome; association only, and
  explicitly not causal, since both are consequences of one intervention.
\end{description}

Two things must be stated here rather than in Results, because they concern
provenance and not findings. \textbf{First, P4 is null} and we report it as
the failure of the thesis it was written to test. \textbf{Second, the
A-versus-B contrast that this paper is organized around was pre-registered as
secondary and exploratory}, not primary. It is labeled post-hoc everywhere it
appears, and its importance does not change its provenance. Additionally, the
pre-registered P1 manipulation objective --- that enforcement would drive
Planner computation-abstention to ${\approx}1.0$ --- \textbf{failed}, because
the check named a clause the intervention does not enforce; the manipulation
check was re-designated to answer authorship and the failure is reported as
such in \S\ref{subsec:structnotcog}.

Exclusions were fixed in advance: trajectories that fail to parse are counted
as incorrect rather than dropped; in Arm~C the regenerated message is the one
analyzed and the rejection count is reported; no arm is dropped after results
are seen, and an arm that fails its manipulation check is reported as having
failed it.

\section{Experiments}
\label{sec:experiments}

\subsection{Tasks and domains}

\TODO{state the two domains (competition mathematics; biology multiple choice),
the problem-selection rule, and the number of problems per cell actually run}

All arms run on the same matched problem set, so every arm comparison is
paired at the problem level.

\subsection{Actor models}

\TODO{name the three actor models and their sizes, the serving stack, and the
pinned backend per actor; report the fraction of trajectories on each backend}

\subsection{Scale and concurrency}

\TODO{report n/cell, total trajectories, total actor calls, total tokens, and
measured wall-clock}

Concurrency was set per actor at validated levels rather than assumed ones
(\TODO{state the per-actor concurrency used and the throughput measured on the
live endpoint}). The panel's largest model is substantially faster per call
than the mid-sized ones, so arms do not complete in equal wall-clock at equal
concurrency.

\subsection{Power}

Power was computed by simulation (4{,}000 replicates, two-sided
$\alpha = 0.05$) from discordance rates and net asymmetries observed in the
existing corpus. At $n=400$ per cell the paired per-actor main effect has power
$0.96$ for large effects but only $0.49$ for moderate ones; detecting a
moderate effect at $0.8$ power requires $n \approx 800$. The
actor~$\times$~intervention interaction (P4), which is the primary test,
requires roughly twice the $n$ of a main effect for equal power, so
\textbf{at the scale run it is powered only for large interactions}. A null on
P4 must be read as ``no large interaction detected,'' not as ``no
interaction.'' \TODO{substitute the achieved n and recompute}

Note the asymmetry: fidelity is measured per \emph{message}, so the measurement
half of the study is well powered at trajectory counts at which the outcome
half is not.

\subsection{Smoke test and launch gates}

A smoke test preceded the main run, and the main run was gated on it. The
gates were: Arm~C produces a non-zero count of rejection/regeneration events
with reasons logged; Arm~C's \Pone{} manipulation check reaches $\approx 1.0$;
answer-bearing calls are equal across arms; \texttt{evaluator\_calls} is
logged; the concurrency ladder is re-measured on the live endpoint; and the
backend is pinned and logged per trajectory.

\TODO{report the smoke-test outcome for each of the six gates, including the
observed rejection count -- this is the evidence that the constraint bound}

\section{Results}
\label{sec:results}

All results below are paired at the problem level with bootstrap clustering by
problem, and P1--P5 carry a Holm correction. Secondary analyses are labeled
exploratory.

\subsection{Did the constraint bind? (manipulation check)}
\label{subsec:manip}

\TODO{report the Arm C rejection/regeneration count per trajectory, the
distribution of rejection reasons, and \Pone{} in Arm C per actor. If \Pone{}
does not reach ${\approx}1.0$, the arm did not engage and every result in
\S\ref{subsec:outcome} onward is void -- report that as the finding.}

\TODO{report answer-bearing calls per arm (the non-handicap check) and
evaluator\_calls coverage}

\subsection{P1: does enforcement move realized fidelity?}

\TODO{table: \Pone{} and \Etwo{} by arm (A, B, C) and actor, with paired
bootstrap CIs on C$-$A}

\subsection{P2: is the fidelity change actor-dependent?}

\TODO{difference-in-differences across actor pairs on fidelity, with bootstrap
CIs}

\subsection{P3: does enforcement move task outcome?}
\label{subsec:outcome}

\TODO{table: accuracy by arm and actor; McNemar on discordant problems for
C vs.\ A, per actor}

\subsection{P4: actor $\times$ intervention on outcome (primary test)}

\TODO{the interaction estimate with CI. This is the paper's primary endpoint.
State the pre-declared interpretation row of Table~\ref{tab:interpretations}
that the observed pattern falls into, before interpreting further.}

\subsection{Capability versus realized fidelity (the discriminating contrast)}

\TODO{report whether the benefit of enforcement is ordered by realized fidelity
or by baseline capability across the three actors. These orderings differ in
this panel, so the comparison is decidable. If it tracks capability, state
plainly that this replicates the capability-equalizer result
\cite{canedo2026architecture} in a multi-agent setting.}

\subsection{P5: does the fidelity change track the outcome change?}

\TODO{per-actor association between $\Delta$fidelity and $\Delta$accuracy
across problems. Association only; not causal -- both are consequences of the
same intervention.}

\subsection{Exploratory analyses}

\paragraph{Role-instruction contribution (A$-$B).}
\TODO{A vs.\ B per actor; a replication of a published null, not a
contribution}

\paragraph{Distinguishability per arm.}
\TODO{recompute function-only balanced accuracy within each arm; does
enforcement change what roles are distinguishable \emph{by}?}

\paragraph{Comparability instruments.}
\TODO{influence rate and restate rate per arm and actor, following
\cite{dylan2026equal}}

\paragraph{Per-domain breakdown.}
\TODO{per-domain results; contracts differ by domain, so components are not
pooled}

\section{Discussion}
\label{sec:discussion}

\paragraph{Why organization would be phase-specific.} The two phases
demand different things. Producing a first answer is a \emph{generation}
problem: there is no incumbent candidate, and three agents arriving at the
same place by three routes costs nothing but tokens, since the pipeline
emits one answer and it is as good as the best route that reached it. A
division of labor has nothing to do here that redundancy does not already
do. Repairing a wrong answer is a \emph{differentiation} problem: an
incumbent candidate exists, it is in every agent's context, and the useful
move is to produce something unlike it. gemma-4's role-free team fails
there not for lack of effort --- it spends its whole budget --- but for
lack of difference: three agents that are now the same solver, all
conditioned on the same incumbent answer, reproduce it. A role contract
supplies a reason for one agent's next output not to be a function of the
previous agent's, worth nothing when there is nothing to differ from and a
great deal when there is; enforcement, the same prohibition moved to the
interface, moves the same quantity. The contract is necessary rather than
sufficient. With roles, Llama-3.3's Reviewer declines in $4.88$ of $5.93$
cycles, yet $0.750$ of its first-pass-wrong trajectories never write down a
second distinct answer and only $10$ of $276$ submit again: the role opens
the revision channel, and this solver has little to send through it.

\paragraph{What a role diagram is evidence of.} A diagram is a
hypothesis, and task-outcome evaluation --- the only evaluation most
multi-agent systems receive --- cannot test it: an aggregate score averages
a stage where the organization is idle with one where it does all its work,
dominated by whichever has more mass. Our system lost its entire
organization and moved the headline metric by $0.0338$. Two cheap additions
fix most of it: report performance \emph{conditioned on the first attempt
having failed}, computable from logs any retry-capable system already
writes, and \emph{measure} the realized organization rather than assume the
specified one --- indistinguishable by accuracy, trivially distinguishable
by whether all three agents are writing down answers.

\paragraph{What our own measurement errors imply.} Three instruments here
reported confidently on what they could not see: an exact-duplication
metric returned $0.000$ for actors a model judge found collapsed; a
logging defect reduced one actor's planner to $37$ visible characters,
which we nearly reported as a behavior; and the adjudication built to
check the first leaked the condition it was meant to be blind to. Each
produced a clean, plausible reading while not measuring the thing. A claim
that a system's structure is intact should therefore state what the
instrument would have to see to report otherwise --- and a blinding claim
should be tested on the bytes the judge saw, not the substitution applied
to them.

\section{Limitations}
\label{sec:limitations}

\paragraph{Power on the primary endpoint.} The actor~$\times$~intervention
interaction requires roughly twice the sample of a main effect for equal power,
and at the scale run it is powered for large interactions only
(\S\ref{sec:experiments}). A null on P4 is ``no large interaction detected.''

\paragraph{Scope.} Three actor substrates, two domains, and one organization
(Planner--Executor--Reviewer). Two of the three actors belong to the same model
family, so the panel varies capability and realized fidelity more than it
varies vendor or architecture. Nothing here establishes that the findings hold
for other organizations, other domains, or other model families.

\paragraph{The capability--fidelity dissociation rests on a narrow base.} In
the observational corpus only one actor \emph{pair} has both a capability proxy
and both fidelity components measured. A correlation over two or three points
is not reported and would not be meaningful. The dissociation is a
hypothesis-generating observation that motivates the design, and the
intervention is what tests it.

\paragraph{Measurement blind spots, stated rather than repaired.} \Pone{} is
blind to symbolic results carrying no digits and therefore slightly
over-credits compliance, so true violation rates are marginally higher than
reported. The variant that would have closed this made agreement substantially
worse and was rejected (\S\ref{sec:measurement}); the limitation is documented
instead. \Etwo{} keys on explicit plan reference and cannot detect an Executor
that follows a plan silently, nor -- as the frozen definitions anticipate --
one that names the plan while ignoring it.

\paragraph{Contracts are domain-specific.} The full role contract exists only
in one domain; other domain presets impose a subset of clauses. Fidelity
components are therefore reported separately and never averaged across domains
or clauses. This is a property of the systems studied, not a choice, and it
limits cross-domain comparison of fidelity levels.

\paragraph{Fidelity is not summarizable.} The two components order the actors
differently (\S\ref{subsec:components-disagree}), so ``most faithful actor'' is
not well defined in this panel, and any mediation claim must name its component
in advance. We do not construct a composite.

\paragraph{The observational association is confounded.} The relationship
between Planner compliance and first-pass success survives removing the
stopping-rule leakage but is confined to problems the baseline already solved,
is null on hard problems, and remains confounded because compliance is itself
difficulty-linked. It is reported as association throughout.

\paragraph{Reviewer behavior is only partially characterized.} The Reviewer
clause measure is defined but not yet validated at the scale of the other two.
\TODO{state whether R4 was annotated and reported, or declare the omission}

\paragraph{Infrastructure confounds.} Serving backends have shown large
throughput and accuracy differences for an identical cell, which is why the
backend is pinned and logged per trajectory. \TODO{confirm that no cell
straddled backends in the run as executed}

\section{Conclusion}
\label{sec:conclusion}

Multi-agent roles need not improve what a team knows on its first attempt.
Their value can appear only after failure, by making external feedback
actionable as revision. Deleting the role instructions destroyed the
division of labor in all three actors while first-pass accuracy was
unchanged or improved; the difference emerged only once both teams were told
their answer was wrong, holding the same six-cycle budget and drawing
\emph{more} oracle verdicts without roles. One team revised and the other
resubmitted, and that split became an outcome split --- cleanly in gemma-4
(paired $p=2\times10^{-8}$, intent-to-treat-equivalent) and again in a
second domain. Making the division of labor structurally unavoidable
un-freezes the answer and returns revisability, the same mechanism from the
other direction.

\textbf{Evaluating agent organization therefore means measuring not just
whether a team is correct, but what its organization lets it do once it
learns it is wrong.} That is scoped here to one \PER{} topology, an
oracle-gated loop, and a recovery effect clean in one actor, directional in
a second and absent in a third.


\balance
\bibliographystyle{ACM-Reference-Format}
\bibliography{refs}

\end{document}
